\documentclass[10pt,a4paper]{amsart}
\usepackage[margin=19mm]{geometry}
\usepackage{amsmath,amssymb,mathtools}
\usepackage[scr=rsfs,cal=euler]{mathalfa}
\usepackage{xfrac}
\usepackage[unicode,hidelinks,bookmarks=false]{hyperref}
\newcommand{\ie}{i.e.\ }
\newcommand{\cf}{cf.\ }
\newcommand{\resp}{respectively\ }
\newcommand{\Subsection}{\subsection}
\providecommand{\nobreakdash}{\nobreak\hbox{-}}
\setlength{\emergencystretch}{2em}
\multlinegap0pt
\usepackage{bm}
\usepackage{bbm}
\usepackage{dsfont}
\usepackage{xspace}
\usepackage{cancel}
\usepackage{mathtools}
\usepackage{amsthm}
\makeatletter
\@ifundefined{Theorem}{\newtheorem{Theorem}{Theorem}}{}
\@ifundefined{Lemma}{\newtheorem{Lemma}[Theorem]{Lemma}}{}
\@ifundefined{Proposition}{\newtheorem{Proposition}[Theorem]{Proposition}}{}
\makeatother
\makeatletter
\newcommand{\mR}{\mathbb{R}}                    %
\newcommand{\norm}[1]{\lVert #1 \rVert}         %
\expandafter\ifx\csname tehtratk\endcsname\relax
\newenvironment{tehtratk}%
             {\begin{list}{\arabic{enumi}.}{\usecounter{enumi}%
              \setlength{\labelsep}{0.5em}%
              \settowidth{\labelwidth}{(\arabic{enumi})}%
              \setlength{\leftmargin}{\labelwidth+\labelsep}}}%
             {\end{list}}
\fi
\makeatother
\title[Generalized boundary rigidity and minimal surface transform]{Generalized boundary rigidity and minimal surface transform}
\author{Leonard Busch, Tony Liimatainen, Mikko Salo, Leo Tzou}
\date{Source: arXiv:2510.23366v1 (CC BY 4.0)}
\begin{document}
\maketitle

\noindent\emph{Source and licence.} Generalized boundary rigidity and minimal surface transform. Version arXiv:2510.23366v1 (CC BY 4.0), distributed under the Creative Commons Attribution 4.0 International licence, as recorded in the arXiv licence metadata for this exact version and shown on the abstract page https://arxiv.org/abs/2510.23366v1. Full rights notice: licenses/arxiv-2510.23366-CC-BY-4.0.txt.\par\smallskip

\noindent\textit{Editorial scope.} Counted unit: the Proposition whose source label is \texttt{prop:evstrict} in the article (source0.tex lines 2053--2059, source label \texttt{prop:evstrict}), delivered together with the article's own complete original proof of it (source0.tex lines 2060--2086, one \texttt{proof} environment of 27 lines). No mathematics is written, repaired or re-derived: statement and proof are the article's own text line for line. Required context retained uncounted: eq:minmax (lines 2026--2028); lem:evmon (lines 2033--2039). Every definition, display and label of those lines is retained unchanged. Omitted from this package and not redistributed: 1--2025, 2029--2032, 2040--2052, 2087--2124. Nothing inside a retained excerpt is omitted or altered: every retained body is delivered exactly as the article's lines are written, so no omission receipt of any kind accompanies this package. Citations: the retained excerpts carry no citation command at all, so no bibliography entry of the article is transcribed and no reference is redistributed. Reference adaptation: none was needed and none was made; every reference inside the delivered unit resolves inside it, so no unresolved reference or citation is produced. Printed slips kept literal: no printed word, symbol, hypothesis or number of the counted statement or of its proof was altered. Layout and class: the article's own class is \texttt{amsart} with packages including amsmath, amscd, amssymb, amsthm, graphicx, mathrsfs, hyperref, tikz-cd, calc, comment, mathtools, xcolor; the wrapper is the lane's pinned \texttt{amsart} editorial wrapper at 10pt on A4, which restates the theorem-like environments the retained text needs and adds the article's own macro definitions that the retained text needs (\texttt{\textbackslash mR}, \texttt{\textbackslash norm}), with extra wrapper packages (bm, bbm, dsfont, xspace, cancel, mathtools). The delivered numbering is the unit's own. Assets: no retained span contains a figure environment, a graphics-inclusion command, a PDF-page inclusion, a table or a TikZ picture; no image file is copied into this package, the package declares no asset dependency and no image file is redistributed with it. Licence: version arXiv:2510.23366v1 of this article is distributed under the Creative Commons Attribution 4.0 International licence, as recorded in the arXiv licence metadata for this exact version and shown on the abstract page https://arxiv.org/abs/2510.23366v1. A full rights notice is delivered at licenses/arxiv-2510.23366-CC-BY-4.0.txt. Characters: every retained excerpt is pure ASCII, so the delivered engine, XeLaTeX with its default Latin Modern text font, reports no missing character.
\begin{equation}\label{eq:minmax}
    \lambda_k(\Omega) = \min_{\substack{\mathcal{L}\subset H^1_0(\Omega)\\\dim \mathcal{L}=k}} \max_{u\in \mathcal{L}\setminus 0} \frac{Q(u,u)}{\norm{u}_{L^2(\Omega)}^2} = \min_{\substack{u\in H^1_0(\Omega)\setminus 0\\ u\perp \mathcal{L}_{k-1}}}\frac{Q(u,u)}{\norm{u}_{L^2(\Omega)}^2}
\end{equation}

\begin{Lemma}\label{lem:evmon}
    For any open $\Omega, \Omega' \subset N$ with $C^1$ boundary, denoting by $\lambda_k(\Omega),\lambda_k(\Omega')$ the $k$-th Dirichlet eigenvalue of $L$ on $H^2(\Omega)\cap H^1_0(\Omega)$ and $H^2(\Omega')\cap H^1_0(\Omega')$ respectively, if $\Omega'\subset \Omega$, then
    \[
        \lambda_k(\Omega)\leq \lambda_k(\Omega')
    \]
    for all $k\in \mathbb{N}$. If there is equality above for all $k$, then $\Omega = \Omega'$.
\end{Lemma}

\begin{Proposition}\label{prop:evstrict}
    Let $\Omega,\Omega'\subset N$ be open with $\Omega\setminus \Omega'$ containing an open set $\subset N$. We have
    \[
        \lambda_k(\Omega) < \lambda_k(\Omega')
    \]
    for all $k\in\mathbb{N}$.
\end{Proposition}

\begin{proof}
    Assume for contradiction that there is $k\in \mathbb{N}$ so that $\lambda \coloneqq \lambda_k(\Omega) = \lambda_k(\Omega')$. Because we know that $\lambda_j(\Omega) \to \infty$ as $j\to\infty$, there must be $m\in\mathbb{N}$ so that $\lambda_m(\Omega)> \lambda$. 

    Define open $\Omega_1,\dots, \Omega_m \subset N$ so that
    \[
        \Omega' = \Omega_1 \subsetneq \Omega_2 \subsetneq \dots \subsetneq \Omega_m = \Omega
    \]
    and $\Omega_{j+1}\setminus \Omega_j$ contains an open set for all $j\in 1,\dots,m-1$. By Lemma~\ref{lem:evmon} we must have
    \[
        \lambda = \lambda_k(\Omega) = \lambda_k(\Omega_m) \leq \lambda_k(\Omega_{m-1}) \leq \dots\leq \lambda_k(\Omega_1) = \lambda_k(\Omega') = \lambda
    \]
    and thus equality in each of these inequalities. 

    Choose eigenfunctions $u_j \in H^2(\Omega_j)\cap H^1_0(\Omega_j)$ of $L$ restricted to the domain $H^2(\Omega_j)\cap H^1_0(\Omega_j)$ with eigenvalue $\lambda$. Denote by $u_j$ the extension of $u_j$ by $0$ to $\Omega$. %
    In order to show that $u_1,\dots,u_m$ are linearly independent in $\Omega$, we consider, for some $a_1,\dots,a_m\in \mR$, $h \coloneqq \sum_{j=1}^m a_j u_j \in H^1_0(\Omega)$. Let us assume that $h=0$. Because $h = a_m u_m$ in $\Omega_m \setminus \Omega_{m-1}$, and because by the unique continuation principle $u_m$ cannot vanish on $\Omega_m\setminus \Omega_{m-1}$, we must have $a_m = 0$. Proceeding in this manner, we find that $a_m=\dots=a_1=0$, so that we have indeed shown that $u_1,\dots,u_m$ are linearly independent. %
    
    Next we consider the function $\varphi := c_1 u_1 + \ldots + c_m u_m \in H^1_0(\Omega)$, where the constants $c_j$ are chosen so that $\varphi$ is $L^2$-orthogonal to the first $m-1$ eigenfunctions of $L$ on $H^2(\Omega) \cap H^1_0(\Omega)$ and $\norm{\varphi}_{L^2(\Omega)} = 1$. %
    By the min-max principle \eqref{eq:minmax}, we then know that
    \[
        \lambda < \lambda_m(\Omega) \leq Q(\varphi,\varphi)\,.
    \]
    On the other hand, using the fact that $u_j$ are $H^2(\Omega_j) \cap H^1_0(\Omega_j)$ eigenfunctions in $\Omega_j$ gives  
    \[
        Q(\varphi,\varphi) = \sum_{j,l} c_j c_l Q(u_j, u_l) = \sum_{j,l} c_j c_l \lambda \langle u_j, u_l \rangle_{L^2} = \lambda \norm{\varphi}_{L^2}^2 = \lambda
    \]
    which gives a contradiction, completing the proof.
\end{proof}

\end{document}
